\subsection{2026-09-22 --- Long format, and two topology bugs}
\label{sec:longform}

\subsubsection*{What we did}

Every pointwise model in this project has been \emph{wide}: one row per (case, hour) with 14
outputs under \texttt{multi\_strategy="multi\_output\_tree"}. That shape was never chosen on
evidence, and it forces three compromises simultaneously:

\begin{enumerate}
  \item \textbf{Shared splits.} A vector-leaf tree picks one split per node and applies it to all
    14 outputs, so the model cannot specialise. Never measured.
  \item \textbf{A 14x-wider count block than necessary.} Forwarding 6 numbers per unit costs 84
    columns (\texttt{cnt\_p\_on\_Lio\_G1}, \texttt{cnt\_p\_on\_Songa\_G1}, \ldots). For Lio's
    head to use Lio's own count, the tree must split on Lio's column --- and that split then
    applies to the other 13 heads, where it is noise. Section~\ref{sec:diagnosis} measured the
    symptom; Section~\ref{sec:abcd} showed the obvious fix was wrong.
  \item \textbf{No per-unit weighting.} \texttt{sample\_weight} applies to a ROW, and a wide row
    carries all 14 units, so a unit cannot be weighted at all.
\end{enumerate}

\texttt{src/longform.py} produces one row per (case, hour, unit) --- 6.9M rows at full stride.
Micro-average ROC-AUC pools all (case, hour, unit) cells into one ranking, so in long format the
metric is simply \texttt{roc\_auc\_score(y, p)}: \textbf{the objective is native to this shape}
and awkward in the wide one.

\subsubsection*{Results}

\begin{center}\small
\begin{tabular}{llrrr}
\toprule
arm & configuration & feats & micro-AUC & Lio \\
\midrule
--- & wide best (Section~\ref{sec:temporal}) & 135 & 0.9442 & 0.8190 \\
\midrule
L1 & long, 15 shared $+$ 6 \texttt{own\_}, no unit id & 29  & 0.9396 & 0.8102 \\
L2 & $+$ unit identity (categorical)                  & 30  & 0.9399 & 0.8090 \\
L3 & $+$ cross-unit context                           & 33  & 0.9421 & 0.8272 \\
L4 & $+$ hourly structure                             & 69  & 0.9456 & 0.8276 \\
L5 & $+$ top-120 count features per unit              & 189 & 0.9490 & 0.8337 \\
L6 & $+$ Lio rows weighted 3x                         & 189 & 0.9479 & 0.8297 \\
L7 & L5 at full stride (5.2M rows)                    & 189 & \textbf{0.9499} & \textbf{0.8380} \\
\bottomrule
\end{tabular}
\end{center}

L1--L6 train on every third case; only L7 uses all of them.

\begin{finding}
Long format reaches \textbf{0.9499}, against 0.9442 for the best wide model. The shape alone is
not the win --- L1 is \emph{below} the wide baseline. The win comes from what the shape makes
affordable: cross-unit context ($+0.0022$), hourly structure ($+0.0035$) and per-unit count
features ($+0.0034$), none of which fit in the wide layout at acceptable cost.
\end{finding}

\begin{finding}
\textbf{Count features at the hour level work in long format and failed in wide format}
(Section~\ref{sec:orthogonal}: 0.9292 raw-forwarded vs 0.9402 distilled). The idea was sound; the
wide layout made each feature cost 14 columns, so 120 features meant 1{,}680 columns of mostly
noise. In long format each costs one column and lands on the row it belongs to.
\end{finding}

\begin{finding}
\textbf{Per-unit weighting fails even where it is finally possible.} L6 weights Lio's rows 3x and
loses $-0.0011$ overall \emph{and} $-0.0040$ on Lio itself. Third independent rejection of the
reweighting family, after historical marginals (Section~\ref{sec:diagnosis}) and per-unit
calibration (Section~\ref{sec:temporal}).
\end{finding}

\subsubsection*{Per-unit: the broken units finally move}

\begin{center}\small
\begin{tabular}{lrrr}
\toprule
unit & wide base & L7 & $\Delta$ \\
\midrule
Haukeli\_G1 & 0.8789 & \textbf{0.9082} & $+0.029$ \\
Lio\_G1     & 0.8157 & \textbf{0.8380} & $+0.022$ \\
Hogga\_G1   & 0.8878 & \textbf{0.9076} & $+0.020$ \\
Vinje\_G1   & 0.9517 & 0.9661 & $+0.014$ \\
Songa\_G1   & 0.9350 & \textbf{0.9059} & $\mathbf{-0.029}$ \\
\bottomrule
\end{tabular}
\end{center}

Three of the four chronically weak units move substantially, for the first time in the project.
Songa moves the other way, and that turned out to be diagnostic.

\subsubsection*{Two topology bugs, found by asking why Songa regressed}

In wide format Songa's head could lean on other units' columns. In long format its row carries
\texttt{own\_*} features and nothing else --- so if those are missing, Songa has nothing. They
are missing.

\begin{finding}
\texttt{plant\_graph} matched only DIRECT reservoir$\to$plant edges, but Songa is fed through
tunnels (\texttt{Songavatn} $\to$ \texttt{Songavatn\_Songa} $\to$ \texttt{Songa}, and likewise
from \texttt{Bitdalsvatn}). So \texttt{src["Songa"]} was \texttt{None}, every \texttt{own\_*}
feature for Songa was NaN, and \texttt{mwh\_per\_m3("Songa", .)} returned NaN --- while Songavatn
carries the \textbf{largest storage variance of any reservoir in the system} (sd 150) and
Bitdalsvatn adds another (sd 25). Section~\ref{sec:v2rejected} diagnosed Songa as ``tunnel-fed''
and attacked it with water-balance corrections; the defect was in the graph traversal.
\end{finding}

\begin{finding}
Three reservoir-to-reservoir TUNNEL links were invisible to the same traversal:
\texttt{Totak}$\to$\texttt{Vaamarvatn}, \texttt{Bordalsvatn}$\to$\texttt{Foersvatn},
\texttt{Venemo}$\to$\texttt{Totak}. The first breaks the main cascade: Songa discharges into
Totak, which tunnels into Vaamarvatn, which is Vinje's source. \textbf{Vinje's upstream plant set
was empty and should contain Songa.} The second hides Bordalsvatn (sd 43) from Kjela.
\end{finding}

Fixed in \texttt{production\_features}: \texttt{\_through\_tunnels}, \texttt{plant\_sources},
\texttt{reservoir\_tunnels} and \texttt{effective\_upstream\_plants}. \texttt{plant\_graph} keeps
its single-source contract for existing callers but now resolves through tunnels, taking the
largest-capacity reservoir as primary where a plant has several.

\subsubsection*{A third finding, from the same question}

Checking whether other units share Songa's problem turned up a different one:

\begin{center}\small
\begin{tabular}{llrrr}
\toprule
reservoir & plant & distinct values & sd & cv \\
\midrule
Vatjern     & \textbf{Haukeli} & \textbf{1} & 0.000 & 0.000 \\
Hyljelihyl  & ---              & \textbf{1} & 0.000 & 0.000 \\
Bandak      & \textbf{Hogga}   & 136        & 8.961 & 0.068 \\
Vinjevatn   & Tokke            & 201        & 1.134 & 0.177 \\
Byrtevatn   & Lio              & 1369       & 20.118 & 0.367 \\
Foersvatn   & Kjela            & 2265       & 33.618 & 0.470 \\
\bottomrule
\end{tabular}
\end{center}

\begin{finding}
\textbf{Haukeli's source reservoir is a literal constant} --- 0.150 for all 3{,}654 days. The
post-decision storage trace, the single largest source of signal in the project
(Section~\ref{sec:production}), \emph{does not exist} for Haukeli. Hogga's Bandak has 136 distinct
values in ten years. The ``four broken units'' were never one problem: Songa was a graph bug,
Haukeli and Hogga are missing data, and \textbf{Lio is neither} --- Byrtevatn varies normally, so
Lio's difficulty remains unexplained.
\end{finding}

\subsubsection*{What this implies}

The standing figure is \textbf{0.9499} on validation year 2021, from 0.9402 at the start of the
session. The topology fix is re-running: it predicts specifically that Songa should recover from
0.9059, since it gains two real reservoirs where it previously had NaN. That prediction is
falsifiable and worth recording before the result arrives.

Caveat: the fix changes features feeding \emph{both} stages, and every experiment to date fitted
the count model around Songa's NaNs. The net effect could go either way.

Still open: Lio, unexplained and worth $+0.0146$ alone; whether any of this survives walk-forward;
and the submission pipeline, which has never been built.
